\section{Coding for Sustaining}\label{sec:CodingForSustaining}
And what to do with \bdq{legacy} code that has to be fixed, amended or extended? 

Most probably that code will not follow \textit{these} coding conventions (if any at all~…), and it is rarely a good idea to reformat or completely rewrite the existing source code just to align it with the current recommendations and guidelines. Instead a fix should be limited to that area in the source code that is broken.

But if possible, your fix should incorporate the standards defined by this document. Be creative how to implement these standards. So if a method has to be reimplemented completely, format it as described here and add the appropriate Javadoc comments\footnote{Even when there are no other Javadoc comments in the whole file~…}, if missing. Or if a new field has to be added to a class, add the ``\verb#m_#'' prefix to its name (see chapter \tqfullref{sec:Fields}).

Obviously, completely new interfaces and classes should be implemented fully compliant to this conventions.

In this context, see also chapter \tqfullvref{sec:MaintenanceComments}.

%==============================================================================
% End of file!!
%==============================================================================
